\section{Experiments}
\subsection{Datasets and Protocol}
We evaluate fringe-break restoration under four-step phase shifting.
DL-SLP contains more than 10,000 physically acquired fringe--height pairs,
whereas SynthFringe (SFNet) maps two fringes to absolute phase. Middlebury
provides stereo images and structured-light-derived disparity, but no projected
fringe sequence. We therefore use each dataset only under its native protocol
and never reinterpret RGB views as phase shifts. Four-step samples are split by
object (70/15/15\%) to prevent leakage.

For each clean sequence, a deterministic degradation generator creates
irregular missing regions, saturated highlights, rectangular occlusions,
Gaussian noise, and salt-and-pepper noise. Missing-area targets are
$r\in\{5,10,20,30\}\%$. The same mask is retained for masked evaluation.

\subsection{Compared Methods and Metrics}
We compare damaged-input PSP, Telea--PSP, Navier--Stokes--PSP, U-Net,
Pix2Pix, direct fringe-to-phase regression, and our dual-head
physics-informed phase restoration network. Image quality is measured by PSNR
and SSIM; circular wrapped-phase error is
\begin{equation}
e_\phi=\operatorname{atan2}(\sin(\hat\phi-\phi),\cos(\hat\phi-\phi)).
\end{equation}
We report phase RMSE, depth RMSE in calibrated metric units, symmetric Chamfer
distance, parameter count, FLOPs, latency, and FPS.

\subsection{Objective}
Given $I_k=A+B\cos(\phi+k\pi/2)$, the total objective is
\begin{equation}
\mathcal L=\mathcal L_{\rm phase}+
\lambda_s\mathcal L_{\rm shift}+
\lambda_g\mathcal L_{\rm grad}+
\lambda_d\mathcal L_{\rm depth}+
\lambda_f\mathcal L_{\rm fringe}.
\end{equation}
The phase term uses circular distance $1-\cos(\hat\phi-\phi)$, avoiding the
$-\pi/\pi$ branch-cut artifact. The shift term reprojects the predicted phase
to all four observations. The gradient term emphasizes damaged boundaries,
and the geometry term is evaluated only at valid calibrated pixels.

\begin{table}[t]
\centering
\caption{Main comparison (mean$\pm$std over three seeds). Dashes must not be
replaced by uncalibrated proxy values.}
\begin{tabular}{lccccc}
\hline
Method & PSNR$\uparrow$ & SSIM$\uparrow$ & Phase RMSE$\downarrow$
& Depth RMSE$\downarrow$ & FPS$\uparrow$\\ \hline
PSP & -- & -- & -- & -- & --\\
Telea & -- & -- & -- & -- & --\\
U-Net & -- & -- & -- & -- & --\\
Pix2Pix & -- & -- & -- & -- & --\\
Fringe-to-Phase & -- & -- & -- & -- & --\\
Ours & -- & -- & -- & -- & --\\ \hline
\end{tabular}
\end{table}

\subsection{Ablations and Discussion}
We remove the shift, phase, and depth terms independently, test all four damage
ratios, and train/test across datasets. High PSNR need not imply low 3-D error:
pixel losses favor smooth intensity interpolation, while phase is encoded by
the relative quadrature among frames; a small correlated intensity bias can
produce a large angular error. Reprojection constrains this inter-frame
relation and reduces physically impossible solutions. Robustness claims are
made only from the ratio sweep, and cross-dataset claims only from
object-disjoint train--test transfers. Industrial deployment additionally
requires radiometric linearization, projector-camera calibration, uncertainty
masking, and tests on specular and discontinuous parts.
